\documentclass[11pt,a4paper]{article}
\usepackage[margin=25mm]{geometry}
\usepackage{amsmath,amssymb,amsthm,hyperref}
\newtheorem{theorem}{Theorem}
\newtheorem{lemma}[theorem]{Lemma}
\newtheorem{corollary}[theorem]{Corollary}
\newcommand{\gen}[1]{\langle#1\rangle}
\title{Paired orbit cuts in string actions}
\author{Human Compression project: incomplete L044 common supplement}
\date{}
\begin{document}\maketitle
This working section rewrites the elementary fracture exchanges in
Cameron--Fernandes--Leemans--Mixer, \emph{Highest rank of a polytope for
$A_n$}, arXiv:1605.09173v1, Section~4. It is not a proof of that paper's
main theorem. The orbit-cut formulation below separates statements that
hold for every string action from the later alternating-group arguments.
The global middle-colour exchange in Lemma~\ref{global} is a new
proof device supplied by the root collaborator; it replaces the local
detour in the source's square-normalisation argument. The theorem and
fracture-graph framework remain attributed to the source. This does
not constitute a new solution of its already proved rank theorem.

Let $G=\gen{\rho_0,\ldots,\rho_{r-1}}$ act faithfully and transitively
on a finite set $\Omega$. Each $\rho_i$ is an involution, and
$\rho_i\rho_j=\rho_j\rho_i$ for $|i-j|>1$. Put
$G_i=\gen{\rho_j:j\ne i}$. The permutation graph has an $i$-edge
$\{x,x^{\rho_i}\}$ for every nonfixed pair. Distinct colours are retained
as distinct edges even when they have the same endpoints. Call an
$i$-edge \emph{eligible} if its endpoints lie in different $G_i$-orbits.
A paired fracture selection $Q$ chooses exactly two eligible edges of
each colour, on the full vertex set. It need not be connected.

\begin{lemma}[Orbit cuts and exchange]\label{cut}
An eligible $i$-edge cannot lie on a cycle whose other edges all have
colours different from $i$. Every simple cycle of $Q$ contains either
zero or two edges of each colour. In particular $Q$ has no parallel
edges. If a cycle of the full permutation graph contains exactly two
$i$-edges and one is eligible, the other is eligible too. Consequently
either of the selected $i$-edges can be replaced by any unselected
$i$-edge so obtained.
\end{lemma}
\begin{proof}
Paths avoiding colour $i$ lie within $G_i$-orbits. This proves the
first assertion and excludes a cycle containing exactly one selected
$i$-edge. There are only two selected edges of that colour, proving
the second assertion. Two parallel selected edges would have different
colours and form a forbidden two-edge cycle. For the last assertion,
delete the two $i$-edges from the cycle. Its two remaining paths show
that the endpoints of the second edge belong to the two distinct
$G_i$-orbits of the first edge. Replacing a selected edge preserves
exactly two eligible edges of the colour and leaves all other colours
unchanged.
\end{proof}

\begin{lemma}[A selected square is a block]\label{squareblock}
An alternating square of $Q$ shares no edge with another simple cycle
of $Q$. More strongly, no path whose internal vertices avoid the square
can join two distinct square vertices using no square edges.
\end{lemma}
\begin{proof}
Write the square as $v_0v_1v_2v_3v_0$, with successive colours $i,j,i,j$.
All selected edges of colours $i,j$ already lie in this square. An
external path between two distinct square vertices uses neither colour.
If the endpoints are adjacent, append the single square edge between
them. If they are opposite, append a two-edge square arc. In either
case the resulting simple cycle contains exactly one $i$-edge or
exactly one $j$-edge, contradicting Lemma~\ref{cut}. Any other cycle
sharing a square edge would have such an external segment.
\end{proof}

\begin{lemma}[Long-cycle reduction]\label{long}
If a paired fracture selection exists, one exists in which every simple
cycle is an alternating square. Neither evenness, primitivity nor the
C-group intersection property is required.
\end{lemma}
\begin{proof}
Take a simple cycle $C$ with length greater than four. Let $i$ be its
smallest colour. The two $i$-edges are disjoint because the $i$-edges
form a matching. At least three distinct edges of $C$ meet them: if
only two did, $C$ would itself have length four. Only two selected
edges have colour $i+1$. Thus some pair of adjacent edges of $C$ has
colours $i,j$ with $|i-j|>1$.

The corresponding commuting involutions produce an alternating square
$S$ in the full permutation graph. Its four vertices are distinct: two
adjacent differently coloured edges are distinct, and commutation of
two involutions gives either their four-point orbit or a common
transposition; the latter cannot contain these two adjacent edges.
There is no need to assume that both completing edges of $S$ are
unselected. For each missing completing edge, replace the selected
edge of the same colour that lies outside $S$. Such an edge exists
because exactly two edges of that colour are selected; the edge of
the initial adjacent pair remains selected. All selected edges of
colours $i,j$ lie in $C$ by Lemma~\ref{cut}, so every removed edge
belongs to $C$. Both completing edges are eligible, by applying the
exchange part of Lemma~\ref{cut} to $S$.

At least one completing edge is missing: otherwise $Q$ would already
contain the square $S$ and the distinct cycle $C$ sharing an edge,
contrary to Lemma~\ref{squareblock}. Make the one or two indicated
exchanges. This preserves exactly two eligible edges of each colour
and destroys $C$.
The new selection contains all of $S$. By Lemma~\ref{squareblock},
neither added edge can belong to any other cycle. Every new long cycle
would therefore avoid both added edges and already have been a long
cycle of $Q$. The original $C$ is destroyed. The finite number of
long simple cycles strictly decreases. Repeating gives the result.
\end{proof}

\section*{Moving square blocks}
The following details make explicit the component collisions in the
source's square moves. Throughout this section every cycle of $Q$ is
an alternating square. Thus every edge outside its square blocks is a
bridge. Two different squares have disjoint colour sets, since either
one uses both selected edges of each of its colours.

\begin{lemma}[Merging incident squares]\label{incident}
If two selected squares meet at a vertex, there is a paired selection
having only square cycles and strictly fewer squares.
\end{lemma}
\begin{proof}
Let the squares have colour sets $\{a,b\}$ and $\{c,d\}$ and common
vertex $v$. They have no other common vertex by
Lemma~\ref{squareblock}. Among the four distinct colours choose one
from each square, called $a,c$, with $|a-c|>1$. Such a choice exists
because two distinct integers have at most one common neighbour in
the integer path. Put $z=v^{\rho_a\rho_c}$.
Commutation gives a square with consecutive vertices
$v,v^{\rho_a},z,v^{\rho_c}$. Its two completing edges are eligible
by Lemma~\ref{cut}.

The vertex $z$ is outside the component of $v$ in $Q$. Indeed every
$Q$-path from $v$ to $z$ would have to use colour $a$: otherwise
prepend the $c$-edge from $v^{\rho_c}$ to $v$ and contradict
eligibility of $\{v^{\rho_c},z\}$. The same argument requires
colour $c$. But all selected $a$-edges lie in the first square and
all selected $c$-edges in the second. No simple path starting at
their common cut vertex can enter both blocks. This is a contradiction.

Delete either selected $a$-edge in the first square and add
$\{v^{\rho_c},z\}$. Deleting the square edge does not disconnect
its component; the new edge joins two previously distinct
components. Hence no new cycle is created and precisely the first
square is destroyed. Eligibility and colour multiplicities are preserved.
\end{proof}

\begin{lemma}[Ordinary transport, including a repeated path colour]
\label{ordinary}
Let $P$ be a shortest path in $Q$ between two distinct square blocks
$S,T$, of positive length. Suppose its first edge is $\{v,w\}$ of
colour $\ell$, where $v\in S$, and $\ell$ commutes with some colour
$a$ of $S$. There is a paired selection having only square cycles
and either fewer squares, or the same number of squares with two
squares at distance less than $|P|$.
\end{lemma}
\begin{proof}
Put $v'=v^{\rho_a}$ and $w'=w^{\rho_a}=(v')^{\rho_\ell}$.
The new $a$-edge $\{w,w'\}$ and $\ell$-edge $\{v',w'\}$ are
eligible. Both $w,w'$ lie outside $S$: the $a$-matching already
pairs every vertex of $S$ internally. Delete the $a$-edge of $S$
opposite $\{v,v'\}$ and add $\{w,w'\}$. If the parallel
$\ell$-edge $\{v',w'\}$ is not already selected, add it and
delete the selected $\ell$-edge other than $\{v,w\}$.
This completes the square $S'=vv'w'w$ and destroys $S$.
Each added edge lies in $S'$, so Lemma~\ref{squareblock} shows
that it creates no other cycle. Consequently no long cycle is
created and the number of squares does not increase.

If the deleted second $\ell$-edge is outside $P$, then either a
further square was destroyed, or $T$ and the path from $w$ to $T$
survive. In the latter case $S'$ is closer to $T$ by at least one
edge. The same argument applies if the parallel edge was already
selected and no $\ell$-edge was removed.

It remains to handle the case where the second selected $\ell$-edge
lies farther along $P$. Here use a different exchange.
Let $C$ be the original component and let $W$ be the component on
the $w$ side after deleting the bridge $\{v,w\}$. Both selected
$a$-edges are outside $W$. If $w'\in W$, a path from $w$ to $w'$
avoids $a$, contradicting eligibility of $\{w,w'\}$. If
$w'\in C\setminus W$, a path from $v'$ to $w'$ avoids $\ell$,
since both selected $\ell$-edges lie on or beyond the cut towards
$W$. This contradicts eligibility of $\{v',w'\}$. Thus $w'$
lies in another component of $Q$.

Delete the opposite $a$-edge of $S$ and the \emph{first}
$\ell$-edge $\{v,w\}$, and add $\{w,w'\}$ and $\{v',w'\}$.
The first deletion destroys $S$ without disconnecting $C$; the
second splits $C$ into two pieces. The new edges join these pieces
through the previously separate component containing $w'$. These
two joins form a tree on three components, so no cycle is created.
The total number of squares therefore strictly decreases.
\end{proof}

\begin{lemma}[Global middle-colour exchange]\label{global}
Suppose a selected square $S$ has colours $i-1,i+1$. Form $Q'$ by
retaining the four edges of $S$ and applying the vertex permutation
$\rho_i$ to every other selected edge. Then $Q'$ is a paired fracture
selection with exactly the same number of squares and no other cycles.
If a shortest selected path from $S$ to a different square starts with
an $i$-edge, then in $Q'$ the distance between $S$ and the transported
target square is strictly smaller than the length of that path.
\end{lemma}
\begin{proof}
Every selected edge of colour $i-1$ or $i+1$ already belongs to $S$.
For every other colour $k$, the vertex permutation $\rho_i$ commutes
with $\rho_k$ (including $k=i$). Thus it sends a $k$-edge to another
actual $k$-edge, injectively. If $k\ne i$, then $\rho_i\in G_k$
preserves each $G_k$-orbit, so the image of an eligible edge is eligible.
For $k=i$ the edge itself is preserved setwise. Hence colour counts
and eligibility are unchanged.

As an abstract coloured graph $Q^{\rho_i}$ has exactly the transported
old square cycles. To obtain $Q'$, delete the transported four edges
of $S$ and insert the original four edges of $S$. No other old square
uses either of these colours, so every other transported square
survives. Every inserted edge lies in the complete selected square
$S$ in the valid paired selection $Q'$. By Lemma~\ref{squareblock}
it belongs to no other cycle. It follows that the only cycles are
$S$ and the other transported old squares, proving the count assertion.

Write the first edge of the path as $v\mathrel{-_i}w$ with $v\in S$.
The rest of this path uses neither $i-1$ nor $i+1$, since all selected
edges of those colours lie in $S$ and a shortest path exits that block
only once. Transport the rest of the path by $\rho_i$. Its initial
vertex is $w^{\rho_i}=v\in S$, while its final vertex lies in the
transported target square. All its edges belong to $Q'$. Its length
is one less than the old path. The target remains a distinct square,
since its colour set is disjoint from that of $S$. Vertex collisions
can only shorten the distance; no assumption about new vertices or
nonrepeated connector colours is needed.
\end{proof}

\begin{theorem}[Square normal form]\label{normal}
If a paired fracture selection exists, then one exists whose every
connected component is a tree or is unicyclic with an alternating
square as its unique cycle. Consequently
\[
                            2r\leq |\Omega|.
\]
If equality holds, every component of this selection is unicyclic.
\end{theorem}
\begin{proof}
First apply Lemma~\ref{long}. Among all selections whose only cycles
are squares choose one with the fewest squares. If a component
contains two squares, choose two at minimum distance, and minimise
this distance over all selections with that fewest-square count.
These minima exist because the full permutation graph is finite.

The distance cannot be zero, by Lemma~\ref{incident}. Let the first
edge of a shortest connector have colour $i$. It lies outside any
square and is therefore a bridge. The colour $i$ differs from both
colours of the first square, since that square uses both selected
edges of each of its colours. If $i$ commutes with either colour of
the square, Lemma~\ref{ordinary} contradicts one of the two minima.
Otherwise the square colours are precisely $i-1,i+1$; the global
exchange of Lemma~\ref{global} preserves the number of squares and
decreases the distance, again a contradiction. Hence no component
contains two squares.

A tree on $v$ vertices has $v-1$ edges and a unicyclic component has
$v$ edges. Summing over components gives the inequality, and gives
equality exactly when there are no tree components.
\end{proof}

\begin{corollary}[Exact defect accounting]\label{defect}
For a selection supplied by Lemma~\ref{long}, let $c$ be its number of
connected components and $q$ its number of squares. Then
\[
                    2r=|\Omega|-c+q.
\]
Consequently $q<c$ implies $2r\leq |\Omega|-1$.
\end{corollary}
\begin{proof}
Each cycle is a square, and distinct squares share no edge by
Lemma~\ref{squareblock}. Delete one edge from every square. The result
has no cycles and the same connected components. It has
$|\Omega|-c$ edges, proving the identity.
\end{proof}

\section*{A single crossing edge}
The complementary branch begins with a colour for which only one
eligible edge exists. The following statement isolates what this
hypothesis actually supplies.
\begin{lemma}\label{single}
If colour $i$ has exactly one eligible edge $\{a,b\}$, then $G_i$ has
exactly two orbits, containing $a,b$ respectively. Every generator
$\rho_j$ with $|j-i|>1$ fixes both $a$ and $b$.
\end{lemma}
\begin{proof}
Contract the $G_i$-orbits in the connected full permutation graph.
Only eligible $i$-edges connect distinct vertices of the quotient. A
connected graph with exactly one edge has exactly two vertices.
For the second assertion, $\rho_j\in G_i$ preserves each of these
orbits. Since it commutes with $\rho_i$, it sends an eligible $i$-edge
to an eligible $i$-edge. Uniqueness makes it preserve $\{a,b\}$, and
preserving the two distinct orbits forces it to fix the endpoints
separately.
\end{proof}

\begin{lemma}[No gap in a path of colours]\label{gap}
Under Lemma~\ref{single}, let $i<k<\ell$. Every path from $a$ or $b$
to a vertex moved by $\rho_\ell$ uses colour $k$. Dually, if
$\ell<k<i$, every such path to a vertex moved by $\rho_\ell$ uses
colour $k$.
\end{lemma}
\begin{proof}
Suppose a path omits colour $k$. Its word factors as a word $u$ in
colours less than $k$ followed by a word $v$ in colours greater than
$k$: these two subgroups commute by the string condition. Every
letter of $v$ fixes the starting point by Lemma~\ref{single}.
Commuting $v$ to the beginning shows that the endpoint is the image
under $u$ alone. But $\rho_\ell$ also fixes the starting point and
commutes with every letter of $u$. It therefore fixes the endpoint,
a contradiction. Reverse the inequalities for the dual case.
\end{proof}

\section*{Outstanding uses}
The remaining two-fracture task is to rule out the equality
$2r=|\Omega|$ for a primitive alternating action. Theorem~\ref{normal}
does not yet show that the all-unicyclic equality case can be made
connected, or that it preserves a system of pairs. The source's
connected-ladder and disconnected-merging arguments remain to be
recovered with their full hypotheses. For odd $|\Omega|$, the bound
already improves to $2r\leq |\Omega|-1$ by integrality.
For the one-crossing-edge branch, the path lemma supplies forced
intermediate colours but does not yet prove the fixed-support budget,
the small $A_5$ end cases or the complete interval induction.
All these remain open obligations in this draft. No complete L044
recovery, new conjecture solution or formal theorem is asserted.
\end{document}
